\section{Introduction}

Large language models are trained on massive human language corpora,
which induce a characteristic geometry in learned representations.
Across layers and training stages, activations, gradients, and model parameters
exhibit pronounced anisotropy:
a small number of directions concentrate disproportionate energy
(strong spectral spikes),
while the remaining directions form a broad long tail
\cite{ethayarajh2019contextual,mu2018all,timkey2021rogue,kovaleva2021bertbusters,machina2024anisotropy}.

Although such structure reflects statistical regularities of language,
shared grammatical patterns versus diverse semantic content, it creates computational challenges.
Anisotropy inflates elementwise dynamic range,
distorts adaptive second-moment estimates \cite{kingma2014adam,loshchilov2017adamw,spectra2026},
and compresses long-tail variation in low-bit regimes
\cite{xiao2023smoothquant,frantar2022gptq,lin2024awq,li2024svdquant,ashkboos2024quarot}.
Consequently, training becomes numerically unstable or inefficient
unless spectral structure is explicitly controlled.

In this work, we investigate the origin of anisotropy across activations, gradient outputs, and model parameters. We uncover a structural asymmetry: \emph{anisotropy in activations is dominated by a coherent mean bias, whereas gradient outputs exhibit no comparable mean component}. Tracing this structure backward, we identify its source in the embedding layer. Because token frequencies in natural language are highly skewed, embedding weights develop a frequency-aligned bias \cite{mikolov2013efficient,mikolov2013distributed,arora2017simple}. Through residual connections, this bias propagates across layers, inducing persistent anisotropy throughout the network.

This reframes anisotropy at its root.
Rather than being a diffuse high-dimensional artifact,
it is largely anchored in a structured, frequency-dominated mean component.
Crucially, this structure is not merely harmful—it is also exploitable.
Because the dominant instability arises from a rank-one bias,
it can be isolated and removed with a single mean-subtraction operation,
consistent with prior embedding post-processing approaches
\cite{arora2017simple,mu2018all,huang2021whiteningbert}.
The remaining low-rank spike can, if necessary,
be whitened once at the embedding activation.
Both operations are hardware-friendly and substantially cheaper
than repeated layer-wise SVD corrections.
In this sense, the same structure that creates instability
also enables efficient mitigation—the \emph{blessing} of bias.

We validate this bias-centric anisotrophy in two downstream settings.
In low-bit quantization, mean removal significantly reduces dynamic-range inflation
\cite{xiao2023smoothquant,wei2023outlier,ashkboos2024quarot}.
In adaptive optimization, it contracts extreme second-moment tails,
improving effective learning-rate allocation
\cite{kingma2014adam,spectra2026}.
These results suggest that addressing bias at its source
can replace or complement spectral interventions
with minimal computational overhead.


\paragraph{Contributions.}
\begin{enumerate}[leftmargin=*,itemsep=2pt]
    \item We identify activation mean bias as the primary structural source of anisotropy in large language model training.
    \item We show that mean-dominated components account for most extreme activation magnitudes and second-moment inflation.
    \item We trace this bias to skewed token frequencies in the embedding layer and its propagation through residual connections.
    \item We propose source-level mean removal and spike whitening as efficient spectral conditioning mechanisms.
    \item We demonstrate improvements in low-bit quantization and adaptive optimization.
\end{enumerate}


















\draft{
Large language models are trained on massive human language corpora.
Across architectures and scales, these corpora induce a characteristic geometry in
learned representations.
Empirically, activations and gradients in modern language models are highly
anisotropic: a small number of directions exhibit unusually large energy,
forming strong spectral spikes, while the remaining directions constitute a
broad, information-rich long tail.
This spike--tail structure appears consistently across layers, models, and
training regimes.

This anisotropy is not an artifact of a particular model or optimizer.
Rather, it is a fundamental property of human language data.
Linguistic structure---grammar, syntax, and shared functional patterns---is
highly coherent across tokens and sentences.
In contrast, semantic content is diverse, context-dependent, and weakly
correlated across tokens.
When embedded into a high-dimensional representation space, this statistical
asymmetry manifests as a low-rank linguistic subspace that induces spectral
spikes, alongside a high-dimensional semantic subspace that forms a long tail.

This structure is a \emph{curse}.
Strong anisotropy leads to poorly conditioned optimization landscapes, suppresses
learning in semantically meaningful tail directions, and interacts unfavorably
with adaptive optimizers such as Adam and AdamW.
In low-bit or mixed-precision settings, the dominant spike controls numerical
scale and clipping behavior, causing severe precision loss in the long tail.
As a result, aggressive quantization and large-scale optimization become
unstable or inefficient unless anisotropy is explicitly addressed.

The standard approach to controlling anisotropy is spectral.
Methods based on singular value decomposition (SVD) or related techniques
explicitly identify and suppress the dominant spike subspace, and have proven
effective in both optimization and quantization settings.
However, these approaches are computationally expensive and hardware unfriendly,
requiring QR decompositions or iterative orthogonalization that do not map well
to modern accelerators.
This limits their applicability in large-scale training loops.

In this work, we identify a complementary and previously underexplored aspect of
anisotropy: \emph{anisotropy is not only a curse, but also a blessing}.
The same statistical regularities that give rise to spectral spikes imply that
linguistic structure is coherent and shared across tokens, while semantic
components are diverse and weakly correlated under batch sampling. In other words, linguistic structure is the shared frozen component, while semantics appears primarily as higher-order structured variability captured by covariates.
As a consequence, the dominant spike admits a simple statistical
characterization: a substantial portion of its energy concentrates in shared,
low-dimensional components that can be accessed without explicit spectral
decomposition.

Leveraging this observation, we show that a much simpler operation---\emph{averaging}---can
serve as a principled and efficient spectral approximation.
By subtracting the batch mean, we remove a shared rank-one component that lies
primarily within the low-rank linguistic spike subspace, while preserving the
semantic covariance that encodes learnable information.
Unlike SVD-based spike removal, averaging is conservative: it suppresses the most
harmful shared structure without destroying semantic variation, yet it can
\emph{fill in the role of SVD-based spike suppression} in practice.

We formalize this observation through a simple model that separates linguistic
and semantic components, provide a theoretical analysis showing that averaging
removes shared linguistic structure while preserving semantic covariance up to
vanishing error, and empirically validate the theory using activations from large
language models.

We demonstrate the practical impact of this perspective in two complementary
settings.
First, in low-bit quantization, averaging reduces spike-dominated scale effects
and preserves long-tail resolution, explaining and complementing recent spectral
methods for FP4 training.
Second, in optimization, averaging attenuates spike-dominated second-moment
statistics that suppress tail learning under adaptive optimizers, providing a
lightweight approximation to spike-aware optimization techniques.
In both cases, averaging achieves the desired spectral control while avoiding
costly decompositions.

Taken together, this work highlights a dual role of anisotropy in human language
corpora.
While anisotropy introduces fundamental challenges for optimization and
quantization, it also enables efficient remedies that exploit the coherence of
linguistic structure.
By aligning algorithmic design with data geometry, we obtain a simple,
theoretically grounded, and hardware-friendly approach to mitigating the
pathologies induced by anisotropic language representations.

\paragraph{Contributions.}
This work makes the following contributions:
\begin{itemize}
    \item We characterize anisotropy as a fundamental geometric property of human
    language corpora, arising from the coexistence of shared linguistic structure and
    diverse semantic content, and show how this structure manifests as a low-rank
    spectral spike with a high-dimensional semantic tail in large language model
    representations.
    \item We formalize the dual role of anisotropy as both a \emph{curse} and a
    \emph{blessing}: while the spike--tail structure can degrade optimization and
    quantization by suppressing long-tail directions, the same statistical
    regularities imply that the dominant linguistic component is coherent and shared
    across tokens, enabling efficient isolation.
    \item We develop a theoretical framework explaining why subtracting the batch mean
    removes shared linguistic structure while preserving semantic covariance up to
    vanishing error, clarifying where information relevant to learning resides.
    \item We introduce averaging-based spectral approximation as a principled,
    hardware-friendly alternative that can fill in the role of SVD-based spike
    suppression, and empirically validate its effectiveness on large language model
    activations by explaining and complementing recent advances in low-bit
    quantization and spike-aware optimization.
\end{itemize}
}